\chapter{Morning Sickness (Nausea Gravidarum)}

\section*{Definition}
Nausea and vomiting experienced during pregnancy, typically beginning around week 6, peaking around week 9, and resolving by week 14-16, though symptoms may begin earlier or persist longer.

\section*{What Happens in Your Body}
Rising human chorionic gonadotropin (hCG) and estrogen levels during early pregnancy stimulate the brainstem vomiting center. Other contributing factors include enhanced olfactory sensitivity, delayed gastric emptying, and gastroesophageal reflux. Psychological factors and prior history of motion sickness may increase susceptibility.

\section*{Causes}
\begin{itemize}
  \item Normal pregnancy-related hormonal changes
  \item Elevated hCG levels (more severe in multiple pregnancies or molar pregnancy)
  \item Estrogen surge in first trimester
  \item Enhanced sense of smell
  \item Delayed gastric emptying
  \item Gastroesophageal reflux
  \item Multiparous pregnancy history
  \item Personal or family history of hyperemesis gravidarum
  \item Migraine susceptibility
  \item Psychological factors (stress, anxiety)
\end{itemize}

\section*{When to See a Doctor}
See your doctor if:
\begin{itemize}
  \item Symptoms are severe or getting worse over time
  \item Severe vomiting causing dehydration (inability to keep fluids down), weight loss, or ketosis (hyperemesis gravidarum)
  \item You have other concerning symptoms such as fever, weight loss, or bleeding
\end{itemize}

\section*{Related Symptoms}
\begin{itemize}
  \item Fatigue
  \item Food aversions
  \item Weight loss
  \item Dehydration
  \item Dizziness
\end{itemize}
\textbf{Important:} If this symptom persists, your doctor will check for serious underlying causes.

\section*{Self-Care / Home Management}
\begin{itemize}
  \item Rest and avoid activities that make it worse.
  \item Maintain adequate hydration and a balanced diet.
  \item Over-the-counter remedies may provide symptomatic relief (consult a pharmacist or doctor).
  \item Monitor symptoms and seek professional advice if they persist or worsen.
  \item Avoid self-medicating with prescription medications without medical guidance.
\end{itemize}

\section*{How Your Doctor Will Evaluate You}
\begin{itemize}
  \item A thorough history and physical examination are the first steps in evaluation.
  \item Laboratory tests (blood work, urinalysis) may be ordered based on clinical suspicion.
  \item Imaging studies (X-ray, ultrasound, CT, MRI) may be indicated when structural pathology is suspected.
  \item Specialist referral is arranged when the diagnosis remains uncertain or requires specific expertise.
\end{itemize}

\section*{Treatment Options}
\begin{itemize}
  \item Treatment targets the underlying cause identified through diagnostic evaluation.
  \item Supportive care and symptom management are provided while awaiting definitive therapy.
  \item Medications, lifestyle modifications, and physical therapies may be prescribed as appropriate.
  \item Follow-up ensures treatment effectiveness and allows timely adjustments to the care plan.
\end{itemize}

\clearpage